% Appendix: exam/paper_b_solutions.md.

% The answers' headings are labelled ansb:q:N.
\renewcommand{\answerprefix}{ansb}

\chapter{Model Answers to Paper B}
\label{ansb:ch:answers}

Model answers for Appendix~\ref{exb:ch:paper}, with marks shown per part.

\textbf{Code questions are answered as checklists} of what a correct answer must contain, not as
listings. Two correct answers will not look alike, syntax is not marked, and a worked listing here
would be a solution to the lab that asks for the same thing.

This paper is about failure, so most marks are in the second half of each answer: not what is wrong,
but what it does to a running system that reports nothing. A script that returns a number and a
script that returns the right number are indistinguishable from the outside, and several questions
turn on exactly that.

% ==========================================================================================
\answersection{1}{What leaves the building, and when it does not}{10}

\answerpart{a}{1}
\emph{If I give somebody this binary, what am I obliged to give them along with it?}

Any phrasing that puts the obligation on \textbf{distribution} earns the mark. Answers about
``whether the code is free'' or ``whether we can use it'' do not; using software you have is rarely
the question, and it is distributing it that triggers everything.

\answerpart{b}{3}
\textbf{Wrong: ``all our code.''} Application code that talks to the kernel through system calls is
not a derived work of the kernel, and its licence is your own business.

\textbf{Right:} the \emph{kernel} is GPL-2.0-only, so the kernel you ship, including any
modifications and any in-tree driver you added, must be published; and so must every other GPL
component in the image.

\textbf{The text that settles it} is the syscall exception, which the kernel's \code{COPYING}
declares at its head and \file{LICENSES/exceptions/Linux-syscall-note} states: user programs using
kernel services by normal system calls are not considered derived works. One mark each, and the
third only for naming \code{COPYING} or the syscall exception specifically rather than gesturing at
``the GPL''.

\answerpart{c}{3}
\textbf{Complete corresponding source} for every GPL component in the shipped image: the source at
the exact version shipped, the \code{.config}, the device tree sources, the scripts that control
compilation and installation, and enough instructions to rebuild the binary you shipped. For
\textbf{three years} from the last distribution of that product.

\textbf{The engineering consequence} is that your build must still work in three years' time. That
means archiving the toolchain and not merely naming it, pinning every source and not tracking a
branch, keeping the \code{.config} as a build artefact rather than as somebody's working file, and
being able to identify which source corresponds to which shipped unit. It is a version control and
artefact retention problem, and it lands on whoever owns the build.

One mark for the contents, one for three years, one for a consequence that is about reproducing an
old build rather than about lawyers.

\answerpart{d}{2}
The script counts \textbf{strings, not licences.} SPDX has current and deprecated spellings of the
same licence, and a tree written over several years contains both: \code{GPL-2.0} and
\code{GPL-2.0-only} are one licence, \code{GPL-2.0+} and \code{GPL-2.0-or-later} are another, and
dual-licensed files carry an \code{OR} expression that a naive grep counts as a licence in its own
right.

One mark for identifying the deprecated-alias problem, one for a concrete pair. A candidate who says
only ``the script has a bug'' earns nothing; the question is which bug.

\answerpart{e}{1}
Half a mark each:
\begin{itemize}
  \item \textbf{Nobody labels their own file \code{Proprietary}.} The absence of a tag is what
    proprietary looks like, so the interesting files are the untagged ones and the script does not
    report those at all.
  \item \textbf{The tree is not the product.} The shipped image contains binaries with no source in
    this tree: firmware blobs, prebuilt vendor libraries, the toolchain's own runtime.
\end{itemize}

% ==========================================================================================
\answersection{2}{A number that is not what it says}{12}

\answerpart{a}{2}
Too small, by a factor of \textbf{2.5}.

The figures are in units of \code{USER_HZ}, which is 100 and fixed by the ABI, so the correct
conversion is to divide by 100. Dividing by 250 gives $100/250 = 0.4$ of the true value. One mark
for the factor, one for the direction. A candidate who answers 2.5 without direction earns one.

\answerpart{b}{3}
$447 / 100 = 4.47$ seconds.

It is not a mistake because \textbf{\file{/proc/stat}'s first line sums over all CPUs.} With two
CPUs, 2.50 s of wall clock offers 5.00 s of CPU time, and 4.47 s of it was idle: the machine was
about 89\% idle, which is the sensible reading. One mark for the conversion, two for the
explanation.

A candidate who says ``the machine was idle for 4.47 s'' has read the number correctly and
understood it wrongly.

\answerpart{c}{2}
One mark each:
\begin{itemize}
  \item \textbf{\code{sleep 1} does not sleep for one second}, it sleeps for at least one second,
    and the interval actually measured also contains the script's own fork and exec on both sides.
    Dividing by the requested interval rather than by the elapsed one, which \file{/proc/uptime}
    will give you, builds in an error you never see.
  \item \textbf{The script perturbs what it measures.} Sampling \file{/proc/stat} means forking,
    reading and parsing, all of which cause context switches of their own, and on an otherwise idle
    machine the measurement is a substantial fraction of the activity.
\end{itemize}

\answerpart{d}{2}
A spread that large means the quantity \textbf{has no single value}: context switch
rate on an idle machine is a property of the particular moment, dominated by whatever happened to
run, not a characteristic of the machine. Any single run is a sample, not a measurement.

It should be reported to \textbf{one significant figure at best}, and honestly to none: what you can
report is a range. One mark each. A candidate who says ``take the average of the ten'' has missed
that averaging a quantity with no stable value produces a number with no meaning, and earns one.

\answerpart{e}{3}
Measured over 200 repetitions at \code{HZ=250}, where one jiffy is 4 ms:

\begin{narrowtable}{@{}ll@{}}
  \thead{Call} & \thead{Measured} \\
  \midrule
  \code{msleep(1)} & \textbf{8,000 us} \\
  \code{msleep(4)} & \textbf{8,000 us} \\
  \code{msleep(10)} & 15,998 us \\
\end{narrowtable}

\noindent \textbf{\code{msleep(1)} and \code{msleep(4)} are the same} because
\code{msecs_to_jiffies} rounds up to the same value: both 1 ms and 4 ms become \textbf{one jiffy},
and both then sleep the same two ticks. Asking for a quarter as long buys nothing at all.

\textbf{The extra jiffy comes from \code{schedule_timeout}}, not from \code{msleep}. Its guarantee
is that \emph{at least} the requested number of jiffies pass, and since the call arrives at an
arbitrary point inside the current jiffy, the remainder of that jiffy cannot be counted; so N
requested jiffies cost N+1 ticks. The model $(\texttt{msecs\_to\_jiffies}(n) + 1)$ jiffies predicts
all three exactly.

One mark for the figures, one for the rounding-up collision, one for locating the extra tick in
\code{schedule_timeout}. Figures within a few percent are fine; the two equal ones must be equal.

% ==========================================================================================
\answersection{3}{An option that did not survive}{10}

\answerpart{a}{3}
\code{EXPERT} is off in the \code{defconfig} this build starts from. With \code{EXPERT} off,
\code{PREEMPT_RT}'s \code{depends on} is unsatisfied, so the symbol \textbf{has no prompt and cannot
be set}; the line in the fragment survives the merge into \code{.config} but \code{olddefconfig}
then discards it, because a symbol whose dependencies are unmet takes its default and the default is
off.

Nothing warns, because a fragment line for an unsettable symbol is not an error in the Kconfig
model; it is simply a request that could not be honoured.

Two marks for the dependency chain, one for saying the line is silently dropped rather than
rejected.

\answerpart{b}{2}
\textbf{A build error stops you.} This produces a kernel that builds cleanly, boots, runs, and is
wrong, so it goes on to be measured, compared and reported as an RT kernel. Every number taken from
it is a number about the wrong kernel.

\textbf{The first symptom} is that \file{/proc/version} or \code{uname -a} says \code{PREEMPT} where
it should say \code{PREEMPT_RT}, closely followed by RT features being absent: no \code{[irq/N-...]}
threads for interrupts that should be threaded, and latency figures indistinguishable from the
non-RT kernel. One mark each. Accept ``the latency measurement showed no improvement'' as the
symptom.

\answerpart{c}{2}
\textbf{Re-read the fragment after \code{olddefconfig} and check each of its assignments against the
final \code{.config}}, warning on any that did not survive; the course's own \file{ci/kernel.sh}
does this.

\textbf{What makes it possible} is that the fragments are separate files containing explicit
\code{CONFIG_X=y} assignments, so the \emph{intent} exists in machine-readable form separately from
the result. You cannot perform this check against a hand-edited \code{.config}, because there is
nothing left to compare it to.

One mark for the check, one for identifying the separation of intent from result as the thing that
makes it writable.

\answerpart{d}{3}
Both gaps are \textbf{alignment}, and in neither case did any code disappear.

\emph{34,138 to 41,208.} The object's bytes do not land in one section, and each section that loses
bytes is itself aligned, so it shrinks by whole units rather than by what was removed:

\begin{narrowtable}{@{}lrl@{}}
  \thead{Section} & \thead{Delta} & \thead{Content \code{configs.o} put there} \\
  \midrule
  \code{.rodata} & 36,864 & 34,058, the compressed config and a string \\
  \code{.text} & 4,096 & 76 \\
  \code{.init.text} & 112 &  \\
  \code{.rela.dyn} & 96 &  \\
  \code{.exit.text} & 40 &  \\
  Elsewhere & 0 & 4, the initcall entry \\
  \textbf{Total} & \textbf{41,208} & \textbf{34,138} \\
\end{narrowtable}

\lecturefigure[0.76]{exam/images/ikconfig_sizes.png}
  {A bar chart of the same one-line change measured three ways: 34,138 bytes of content in configs.o,
   41,208 bytes of vmlinux sections, and 65,536 bytes of arm64 Image. The two gaps are annotated as
   page alignment and 64 KiB segment alignment.}
  {fig:ikconfig_sizes}

\noindent Two sections carry almost all of it and both moved by an exact multiple of the 4,096-byte
page: \code{.rodata} by nine pages for 34,058 bytes of content, and \code{.text} by a \textbf{whole
page for 76 bytes}. The remaining 248 bytes are the sections that are not page-aligned. Removing a
thing removes its padding with it, so the section total moves by more than the thing measured.

\emph{41,208 to 65,536.} \code{Image} is a flat binary whose segments arm64 aligns to \textbf{64
KiB}; \code{SEGMENT_ALIGN} is \code{SZ_64K} in \file{arch/arm64/include/asm/memory.h} and
\code{vmlinux.lds.S} aligns each segment to it. So the \code{Image} delta is quantised: it moves in whole 65,536-byte
steps, and here 41,208 bytes of sections happened to cost exactly one. A second option of the same
size need not shrink it again by 65,536, and a small option may not shrink it at all; which it does
depends on where each segment ends relative to the next 64 KiB boundary.

One mark for page alignment of the sections, one for the 64 KiB segment quantisation, one for
concluding that the \code{Image} figure cannot resolve anything below 64 KiB and is therefore the
wrong measurement to quote. \textbf{Link-time garbage collection is the plausible wrong answer} and
earns nothing on its own; nothing was collected, and every byte is accounted for by padding.

Candidates are not expected to reproduce the per-section table from memory. Full marks for naming
both alignments and the direction of each gap.

% ==========================================================================================
\answersection{4}{A module that will not load}{12}

\answerpart{a}{4}
Three causes, one mark each, plus one for a method that actually separates them.
\begin{enumerate}
  \item \textbf{The module that exports the symbol is not loaded.} Nothing provides it.
  \item \textbf{The symbol exists but is not exported.} The provider is loaded, the function is
    there, but it carries no \code{EXPORT_SYMBOL}, so it is invisible to the module loader.
  \item \textbf{The symbol is \code{EXPORT_SYMBOL_GPL} and your module's \code{MODULE_LICENSE} is
    not GPL-compatible.} The loader deliberately does not see it.
\end{enumerate}
\textbf{Distinguishing them:} \code{grep qa_answer /proc/kallsyms}. Absent means cause 1, and
\code{lsmod} confirms. Present means cause 2 or 3, and \code{modinfo -F license} on your own module
settles it: a GPL-compatible licence leaves only cause 2, which the provider's source confirms by
the missing \code{EXPORT_SYMBOL}.

\answerpart{b}{3}
\textbf{Cause 3.} Your module is correct, the provider is loaded, the symbol is present in
\file{/proc/kallsyms}, and the message says \code{Unknown symbol}.

It is misleading because the loader does not report a licence violation; it \textbf{hides GPL-only
symbols from a GPL-incompatible module}, so the lookup genuinely fails and the loader reports,
accurately, that it could not find the symbol. Nothing in the message mentions licensing, and the
symbol is plainly there to anyone who looks.

Two marks for the mechanism of hiding rather than refusing, one for noticing that the message is
technically true. This path is normally reached after a kernel upgrade turns an \code{EXPORT_SYMBOL}
into an \code{EXPORT_SYMBOL_GPL}, since a fresh build would have been stopped at modpost; a
candidate who says so earns full marks comfortably.

\answerpart{c}{2}
\textbf{\code{vermagic} disagreed.} The string built into the \code{.ko} records the kernel version,
SMP, the preemption model, module unloading and \code{modversions}, and it did not match the running
kernel's.

The fix is to rebuild the module against the exact kernel it will load into. \code{modinfo -F
vermagic mine.ko} against the target's own string identifies which field differs; the common cause
in this course is switching between the \code{PREEMPT} and \code{PREEMPT_RT} builds and leaving
stale objects behind. One mark each.

\answerpart{d}{2}
The \textbf{modpost} failure is better.

It arrives \textbf{at build time on your own machine, names the licence, the module and the symbol,
and produces no \code{.ko} at all}, so the defect cannot leave the build. The \code{insmod} failure
arrives on the target, possibly at a customer, says \code{Unknown symbol} and does not mention the
actual cause; diagnosing it requires already knowing this mechanism exists.

One mark for the choice with the build-time argument, one for contrasting the quality of the two
messages.

\answerpart{e}{1}
\code{Disabling lock debugging due to kernel taint}.

\textbf{Lockdep is now off for the rest of the boot.} Every lock-ordering bug that the validator
would have caught passes silently from this point, and the only way back is a reboot without that
module.

% ==========================================================================================
\answersection{5}{A read that returns zero}{12}

\answerpart{a}{8}
Two marks each: one for naming the defect, one for the consequence.

\textbf{\code{return 0} when the FIFO is empty.} Zero means \textbf{end of file}. \code{cat} exits,
a \code{while (read(...) > 0)} loop terminates, and the program concludes the device is finished
when it has merely not produced anything yet. It should block, or return \code{-EAGAIN} under
\code{O_NONBLOCK}.

\textbf{\code{memcpy} to a \code{char __user *}.} No verification that the address belongs to the
caller, no fault handling if the page is not present, and no \code{__user} checking. An unmapped
address oopses inside the driver; a kernel address supplied by the caller is written to happily.
Must be \code{copy_to_user}, whose return is the number of bytes \textbf{not} copied.

\textbf{The copy happens under \code{spin_lock_irqsave}.} \code{copy_to_user} may sleep, because the
destination page can be paged in on demand, and this is atomic context with interrupts disabled.
With \code{CONFIG_DEBUG_ATOMIC_SLEEP} it is a \code{might_sleep} splat; without it, the machine can
deadlock or corrupt. Copy into \code{tmp} under the lock, drop the lock, then copy out.

\textbf{\code{return count} rather than \code{n}.} Claims more bytes were delivered than were
written. Only \code{n} bytes were copied, so the caller takes whatever its own buffer already held
past \code{n} as data from the device: the byte stream is silently corrupted, and a short read,
which is normal, is never seen as one.

Accept, for partial credit within a defect the candidate has otherwise missed, the observations that
\code{priv} is never derived from \code{file->private_data} and that a 64-byte array on a small
kernel stack is unwise. Neither is one of the four.

\answerpart{b}{2}
Half a mark each: the byte count actually delivered, which may be short; sleep until data arrives,
then that count; \code{-EAGAIN}; \code{-ERESTARTSYS}.

\textbf{Zero is correct in none of the four}, and offering it for the second or third is the defect
in part~(a) restated.

\answerpart{c}{2}
With no \code{.poll} method, the VFS applies the default mask, which reports the file as
\textbf{always readable and always writable}. \code{select} therefore returns immediately claiming
the descriptor is ready, the program calls \code{read}, gets \code{-EAGAIN} on its non-blocking
descriptor, calls \code{select} again, and \textbf{spins at 100\% CPU} while appearing to work. On a
blocking descriptor it sleeps in \code{read} instead, and every other descriptor it was watching
waits with it.

The driver must implement \code{.poll}, calling \code{poll_wait} on the same wait queue \code{read}
sleeps on and returning \code{EPOLLIN | EPOLLRDNORM} only when data is actually available. One mark
for the busy loop, one for \code{poll_wait} on the same queue.

% ==========================================================================================
\answersection{6}{A mapping that is wrong}{10}

\answerpart{a}{3}
Nothing is decoded at that address, so the read raises a \textbf{synchronous external abort}. The
kernel prints an oops naming the faulting instruction and the driver's function.

The state left behind: \code{insmod} is killed, but \textbf{the module is half-loaded}, with
whatever it allocated before the read still held and its init function never having returned. It
cannot be removed, because \code{rmmod} requires a module that finished loading, and the resources
are held until reboot. The machine survives; that module slot does not.

One mark for the abort, one for \code{insmod} dying, one for the half-loaded state.

\answerpart{b}{2}
The \textbf{zeros}, decisively.

The abort is immediate, loud, and points at the exact instruction. Reading zeros produces no error
of any kind: the driver believes it has a working mapping, sees an identity register that does not
match and, if it does not check, carries on to configure a device that is not there. The failure
then surfaces somewhere else entirely, at a different time, as a device that does not work, and
nothing in the log connects it to the mapping.

One mark for the choice, one for ``no error is worse than an error''.

\answerpart{c}{2}
\textbf{``It did not crash'' is not evidence of anything.} A wrong mapping has two outcomes and only
one of them is loud, and which one you get depends on what happens to be at the wrong address, not
on how wrong it is. Absence of a fault distinguishes nothing.

The defensive measure that catches both: \textbf{read a known identity register and refuse to probe
unless it holds the expected value.} The abort case never reaches the comparison, and the zeros case
fails it, so one check covers both. One mark each.

\answerpart{d}{3}
Any three, one mark each:
\begin{itemize}
  \item \textbf{The compiler treats it as ordinary memory.} It may cache the value in a register and
    never re-read it, so a loop polling a status register never observes it change; it may merge or
    discard writes it believes are redundant; it may reorder accesses against each other.
  \item \textbf{No ordering.} \code{readl} and \code{writel} carry the barriers that make MMIO
    ordered against other accesses and against DMA. A bare dereference carries none, so writes can
    reach the device in an order the driver never wrote.
  \item \textbf{No \code{__iomem} checking.} Sparse can no longer tell an I/O pointer from a normal
    one, so this class of mistake becomes invisible to the one tool that finds it, throughout the
    driver and not just here.
  \item \textbf{Access width and portability.} The compiler chooses the load width, and many device
    registers require an exact width; and on architectures where I/O needs distinct instructions,
    the code is simply wrong rather than merely fragile.
\end{itemize}

% ==========================================================================================
\answersection{7}{A deadlock that has not happened}{12}

\answerpart{a}{3}
Lockdep does not watch for deadlocks; it \textbf{records the order in which locks are acquired and
builds a graph of those orderings}, one edge per nesting it has ever seen, anywhere, at any time.

Taking A then B adds the edge A to B. Taking B then A later adds B to A, closing a \textbf{cycle}. A
cycle in that graph means an ordering exists under which two threads can deadlock, and lockdep
reports on the edge that closes it, immediately, whether or not anything was concurrent.

Two marks for the graph-of-orderings model, one for saying the report is about the
\emph{possibility} and fires on the first inversion. A candidate who says ``lockdep detected a
deadlock'' has the wrong model and earns at most one.

\answerpart{b}{2}
\textbf{One CPU.} The columns are \textbf{hypothetical}: lockdep constructs the interleaving that
\emph{would} deadlock, to show you what it is worried about. Nothing in the diagram is a record of
anything that happened. One mark each, and the second is the examinable half.

\answerpart{c}{2}
\textbf{Lockdep}, the lock validator, enabled by \textbf{\code{CONFIG_PROVE_LOCKING}}.

The cost is substantial: every acquisition and release is instrumented and checked against the
graph, which is far too slow for production and uses memory that grows with the number of distinct
lock classes. It is also \textbf{switched off entirely by the taint a proprietary module brings}, so
loading one ends validation for the boot. One mark for the name and option, one for a real cost.

\answerpart{d}{3}
Because \textbf{the racing module never took a lock.} Lockdep validates statements about locks that
were acquired; where no lock was acquired there is no event to record, no edge to add, and nothing
to contradict. Losing up to 38\% of the increments involves no incorrect use of any lock, so from
lockdep's point of view nothing happened at all.

\textbf{The general limitation: lockdep finds bugs in the locks you took, not the locks you did not
take.} It cannot know that a variable needed protection, and a clean lockdep report says nothing
whatever about whether your data is protected.

Two marks for the mechanism, one for the limitation stated in general terms.

\answerpart{e}{2}
Both sentences needed:

The ordering only deadlocks when the two paths interleave inside a window that may be a few
instructions wide, and two years of not hitting a narrow window is evidence about the workload, not
about the code; the machine that hits it will be a faster one, a busier one, or one with more CPUs.

And the argument proves too much: it is the same reasoning that says an unlocked counter is fine
because 2,000 increments per thread never lost one, which this course measured going wrong at
20,000.

% ==========================================================================================
\answersection{8}{An interrupt that never stops}{10}

\answerpart{a}{3}
The device's status bit stays set, so the line \textbf{stays asserted}. The GIC sees a level that is
still high, so as soon as the handler returns the interrupt is delivered again, immediately and
forever.

The CPU never leaves interrupt context on that core. \textbf{\code{insmod} never returns}, the
console stops responding, and the machine is unrecoverable without a reset. Two marks for the
level-stays-high mechanism, one for \code{insmod}.

\answerpart{b}{3}
The connection is that \textbf{the CPU is never available to anything else.} RCU needs every CPU to
pass through a quiescent state for a grace period to complete, and a CPU wedged in an interrupt
handler never does. After the stall timeout, the RCU stall detector notices a grace period that has
not ended and reports it.

The NMI is the detector asking the stuck CPU what it is doing, since it cannot be interrupted by
anything less; the backtrace it returns is what actually names the driver. The message mentions
neither interrupts nor your driver because \textbf{RCU is a victim, not the cause}: it is the first
subsystem with a timeout long enough to notice.

Two marks for the quiescent-state argument, one for the NMI's purpose.

\answerpart{c}{2}
It should report \textbf{TIMEOUT}, distinctly from a failure.

A failing test is one that ran, produced a result, and the result was wrong. This produced
\textbf{no result at all}: the guest never reached the end of the test, so nothing is known about
anything else the lecture tests either. Reporting it as a failure would imply the harness knows what
went wrong, and treating it as a normal failure hides that every later test in that boot is
unaccounted for. One mark each.

\answerpart{d}{2}
On an edge-triggered line the controller latches a transition, not a level, so a status bit left set
produces \textbf{no new edge and therefore no further interrupt}. You get exactly one interrupt and
then silence: the FIFO fills, \code{OVERRUN} sets, and the device appears dead while the machine
runs perfectly.

\textbf{Which to ship:} neither is shippable, and the honest answer is the trade. The wedge is
undeniable, fires on the first interrupt, and \textbf{cannot escape testing}. The silent one passes
a short test, passes review, and fails hours later in the field as an intermittent ``device stopped
responding'' with nothing in the log.

The position expected: \textbf{the wedge}, on the ground that a failure which cannot reach a
customer is preferable to one that reliably does. Full marks for an answer arguing the opposite on
recoverability, that a wedged machine in the field is unrecoverable while a dead device may be
restarted, provided it states that trade explicitly rather than merely preferring the quieter
symptom.

% ==========================================================================================
\answersection{9}{A driver that finds its own hardware}{6}

\answerpart{a}{3}
Three quarters of a mark each. Order matters: each step is cheaper than the next and rules out
everything below it.
\begin{enumerate}
  \item \textbf{Does the device exist?} Is the node in the device tree, is its \code{status}
    \code{okay}, and did a platform device appear under \file{/sys/bus/platform/devices/}? No
    device, no probe, and nothing about the driver is relevant.
  \item \textbf{Did the driver register?} Is the module loaded and does
    \code{/sys/bus/platform/drivers/<name>/} exist? A driver that failed to register cannot match.
  \item \textbf{Do the strings match exactly?} The \code{compatible} in the DT against the
    \code{of_match_table} entry, character for character, including the vendor prefix. This is the
    usual answer.
  \item \textbf{Is the match table actually wired in?} \code{.of_match_table} assigned into the
    \code{struct platform_driver}'s \code{.driver} member, and the table terminated by an empty
    entry. A table that exists but is not attached matches nothing and warns about nothing.
\end{enumerate}

\answerpart{b}{2}
\textbf{All three are released at unbind}, while the module stays loaded: the memory region
disappears from \file{/proc/iomem}, the interrupt is freed and its line disappears from
\file{/proc/interrupts}, and the sysfs attributes are removed. Rebinding runs \code{probe} again and
brings them all back.

What it tells you: \textbf{\code{devm_} is attached to the device, not to the module.} The lifetime
it manages is the bind, and unloading the module is merely the most common way for a bind to end,
not the thing being tracked. One mark each.

\answerpart{c}{1}
\textbf{A machine with two of these devices.} The second \code{probe} overwrites the pointer the
first stored, so both devices then operate on the second one's mapping, and the first device
silently stops working while both appear to have probed successfully.

% ==========================================================================================
\answersection{10}{Worse average, better maximum}{6}

\answerpart{a}{2}
A \textbf{throughput} reviewer prefers \code{PREEMPT}: its mean is 449 us against 551 us, so the
average path through the system is 19\% cheaper, and RT's overhead is real work being done on every
interrupt.

A \textbf{real-time requirement} demands \code{PREEMPT_RT}: its maximum is 4,964 us against 11,321
us, a factor of 2.3. A deadline is a statement about the worst case, and no average is evidence
about a maximum.

Both read correctly because \textbf{the table describes a trade rather than an improvement}:
\code{PREEMPT_RT} buys a better bound by paying for it on every sample. One mark for the pair of
readings, one for naming the trade rather than declaring a winner.

\answerpart{b}{2}
The same kernel, measured again, moved its maximum from 11,321 us to 2,502 us: \textbf{a factor of
four and a half, on one kernel, changing nothing.} That is larger than the kernel-to-kernel difference the
table is being used to demonstrate, so \textbf{the table supports no conclusion at all.} The
comparison is inside the noise of its own method.

To draw one safely: many runs of each configuration, reported as a distribution rather than a single
maximum, with the run-to-run spread stated alongside; and for a bound, the interesting statistic is
a high percentile over a long run, not the maximum of two thousand samples. One mark for
invalidating the conclusion, one for what would replace it.

\answerpart{c}{2}
Any two, half a mark each: \textbf{QEMU itself}, whose device emulation and timer delivery sit
directly in the measured path; \textbf{the host kernel}, which schedules the QEMU process and can
preempt it mid-interrupt; the host's own load; the host's power management and frequency scaling.

\textbf{What survives:} the \textbf{direction and the shape} of the comparison. \code{PREEMPT_RT}
trading a worse mean for a better maximum is a real property of what it does, and it is visible here
because the mechanism is real even when the timings are not.

\textbf{What does not survive: every absolute number, and the factor.} Not ``approximately correct''
but worthless, because the dominant term is the host and not the guest. A candidate who carries ``RT
is 2.3 times better'' out of this course has taken the one thing the measurement cannot support. One
mark for the components, one for separating shape from magnitude.
